\chapter{Principios básicos de recuento}
\label{md:recuento}

\section{Suma, producto y correspondencias}

La \textbf{cardinalidad} $|A|$ de un conjunto finito es su número
de elementos. Una biyección entre dos conjuntos permite afirmar
que tienen la misma cardinalidad, traduciendo un recuento a otro
equivalente.

El \textbf{principio de adición} se aplica a casos excluyentes:
si $A_1,\ldots,A_r$ son disjuntos dos a dos, entonces
\[
 |A_1\cup\cdots\cup A_r|=|A_1|+\cdots+|A_r|.
\]
Si los casos se solapan, la suma cuenta elementos varias veces;
hay que corregirla mediante inclusión-exclusión
(capítulo~\ref{md:inclusion-exclusion}).

El \textbf{principio de multiplicación} cuenta elecciones sucesivas.
Si hay $m_1$ posibilidades para el primer paso y, para \emph{cada}
historia de elecciones previas, hay exactamente $m_i$ posibilidades
para el paso $i$, el total es
\[
 m_1m_2\cdots m_r.
\]
Las opciones concretas pueden depender de lo elegido antes;
lo que debe permanecer fijo es su número en cada etapa.
En particular, $|A_1\times\cdots\times A_r|=\prod_i|A_i|$.

Antes de contar conviene precisar qué constituye un resultado,
cuándo se consideran iguales dos resultados, si importa el orden
y si se permiten repeticiones. Elegir primero un modelo correcto
evita aplicar una fórmula a objetos diferentes de los que se buscan.

\section{Principio del palomar}
\label{md:palomar}

Si se distribuyen más de $n$ objetos en $n$ cajas, alguna recibe
al menos dos. Más generalmente, para $n\ge1$ y $p\ge0$ enteros,
\[
 N>np\quad\imp\quad
 \text{alguna caja contiene al menos }p+1\text{ objetos}.
\]
Si todas contuvieran a lo sumo $p$, el total no superaría $np$.
Por tanto, en cualquier distribución de $N$ objetos, alguna caja
tiene al menos $\lceil N/n\rceil$ objetos.

El principio prueba existencia sin localizar necesariamente la caja.
Para aplicarlo hay que definir una asignación inequívoca de objetos
a un número finito de clases.

\begin{ejemplo}
Entre siete enteros, dos tienen el mismo resto módulo $6$:
los seis restos posibles son las cajas. Su diferencia es divisible por $6$.
\end{ejemplo}
